\begin{abstract}
Text-to-video generators still render garbled text on signs, screens and posters. We study \emph{repair}: given a generated clip, a target string and a hint box, make the text read correctly and leave everything else as it was. Generative video text editors produce natural glyphs but keep or invent wrong characters, rewrite text that was already correct, and re-render the background. \method is a training-free, correctness-first pipeline built around three evidence gates: a preservation gate leaves a clip untouched when multi-frame OCR already reads the target; a localization gate lays the new text out on the old glyph instance when the evidence is strong and falls back to the tracked box otherwise; and an appearance gate enables a non-parametric source appearance transport (SAT) that carries glow and material from the old glyphs to the new ones only where it measurably helps. The target is rendered deterministically, so its spelling cannot drift, and only pixels inside the text region change. On 526 generated clips over seven carrier materials in English and Chinese, \method reaches \rOursRead\% readability against \rSteerRead\% for the state-of-the-art editor \steer (\rOursNet{} vs.\ \rSteerNet{} net repairs; a pre-registered second reader confirms the advantage) and keeps the background intact (\rOursBG{} vs.\ \rSteerBG\,dB PSNR outside the text). Raters preferred SAT over plain pasting for material fidelity in 87\% of non-tied judgements. We report, with pre-registered criteria, what did not work---a learned 1.3B refiner, two appearance modules, a glyph-locked material generator and reconstruction without any source glyph---and characterise the appearance gap to generative editing that remains.
\end{abstract}
